\section*{الفصل \ref{ch:b1:water-small-molecules} --- الماء والجزيئات الحيوية الصغيرة}

\begin{solution}{exo:b1:water-small-molecules:1}
الجزيء المنحني \omterm{def:b1:water-small-molecules:water}{القطبي} أكسجينه سالب جزئيًا وهيدروجيناه موجبان
جزئيًا: فهو يحيط الكاتيون بأكسجيناته ويحيط الأنيون بهيدروجيناته
فيستقران (إماهة)، وهو يعقد \omterm{def:b1:water-small-molecules:water}{روابط هيدروجينية} مع الزمر القطبية. أما
الزيت فلا شحنات له ولا زمر قطبية يرتبط بها؛ وجزيئات الماء من حوله
عليها أن تنتظم في قفص، وهذا غير مواتٍ، فيُستبعد الزيت.
\end{solution}

\begin{solution}{exo:b1:water-small-molecules:2}
\omterm{def:b1:water-small-molecules:water}{الرابطة الهيدروجينية} (\qtyrange{4}{20}{kJ/mol}): ازدواج القواعد في DNA،
والبنية الثانوية للبروتين. \omterm{def:b1:water-small-molecules:bonds}{والرابطة الشاردية} (نحو \qty{12}{kJ/mol} في
الماء): الجسور الملحية في البروتينات، وارتباط الركيزة. \omterm{prop:b1:water-small-molecules:properties}{والأثر الكاره
للماء} (ليس رابطة، لكنه مقارب في المجموع): تركيب \omterm{def:b1:membranes-transport:membrane}{الأغشية}، وطي
البروتين. \omterm{def:b1:water-small-molecules:bonds}{وفان دير فالس} (\qtyrange{1}{4}{kJ/mol} لكل منها): الارتصاص
الوثيق لأي سطحين متكاملين، وتوافق الإنزيم مع ركيزته.
\end{solution}

\begin{solution}{exo:b1:water-small-molecules:3}
$\mathrm{pH} = -\log(4\times 10^{-8}) = 7.4$. وعند \omterm{def:b1:water-small-molecules:ph}{pH} 1.5 يكون
$[\mathrm{H^+}] = 10^{-1.5} = \qty{0.032}{mol/L}$.
\end{solution}

\begin{solution}{exo:b1:water-small-molecules:4}
من نحو 5.8 إلى 7.8 (أي p$K_a \pm 1$). وعند 7.4:
$\log(\text{أساس}/\text{حمض})
= 0.6$، أي نسبة $4$.
\end{solution}

\begin{solution}{exo:b1:water-small-molecules:5}
عند \omterm{def:b1:water-small-molecules:ph}{pH} 4.76: \qty{50}{\%}. وعند 7.2:
$\log(\mathrm{A^-}/\mathrm{HA}) =
2.44$، أي نسبة $275$: فيكون
\qty{99.6}{\%} مشرَّدًا. وعند \omterm{def:b1:water-small-molecules:ph}{pH} 2: النسبة $10^{-2.76} =
1/575$: أي
\qty{0.2}{\%} مشرَّد. والصورة الحمضية المتعادلة هي التي تعبر \omterm{def:b1:membranes-transport:membrane}{الطبقة
الدهنية المضاعفة}؛ فيُمتص \omterm{def:b1:water-small-molecules:ph}{حمض} الخل في المعدة، حيث يكون غير مشرَّد،
ويُحتجز أسيتاتًا في \omterm{def:b1:eukaryotic-cell:organelle}{العصارة الخلوية}.
\end{solution}

\begin{solution}{exo:b1:water-small-molecules:6}
\omterm{def:b1:water-small-molecules:ph}{pH} $= 6.8 + \log 1 = 6.8$. وبعد HCl: \omterm{def:b1:water-small-molecules:ph}{الأساس} $40$ \omterm{def:b1:water-small-molecules:ph}{والحمض} $60$:
فتكون \omterm{def:b1:water-small-molecules:ph}{pH} $= 6.8 +
\log(40/60) = 6.62$. وفي الماء النقي يكون
$[\mathrm{H^+}] = \qty{0.01}{mol/L}$: أي \omterm{def:b1:water-small-molecules:ph}{pH} $2$.
\end{solution}

\begin{solution}{exo:b1:water-small-molecules:7}
$500/2.4 = \qty{208}{g}$. ولتبخير جزيء يجب كسر روابطه الهيدروجينية
الأربع كلها؛ و\qty{2.4}{kJ/g} هي \qty{43}{kJ/mol}، أي تقريبًا طاقة
رابطتين لكل جزيء (لأن كل رابطة مشتركة بين اثنين).
\end{solution}

\begin{solution}{exo:b1:water-small-molecules:8}
$\Delta G^{\circ\prime} = -RT\ln 19 = -2.48\times 2.94 =
\qty{-7.3}{kJ/mol}$. وفي \omterm{def:b1:cell-unit-of-life:cell}{الخلية}:
$\Delta G = -7.3 + 2.48\ln(2/0.02) =
-7.3 + 11.4 = \qty{+4.1}{kJ/mol}$:
فيجري التفاعل عكسيًا، نحو غلوكوز-1-فوسفات (كما يجري في تركيب
الغليكوجين).
\end{solution}

\begin{solution}{exo:b1:water-small-molecules:9}
الماء أكثف ما يكون عند \qty{4}{\celsius}؛ والماء الأبرد والجليد أخف
فيبقيان على السطح، فيتكون الجليد فوق ويعزل السائل تحته، فيبقى قرب
\qty{4}{\celsius} طول الشتاء. ولو غاص الجليد لتجمدت البحيرات تجمدًا
تامًا من القاع إلى الأعلى ولما ذابت إلا عند السطح في الصيف؛ ولاستحالت
حياة الماء العذب عبر شتاء.
\end{solution}

\begin{solution}{exo:b1:water-small-molecules:10}
$RT = \qty{2.58}{kJ/mol}$. في الراحة:
$Q = 0.9\times 10^{-3}\times 8\times
10^{-3}/8\times 10^{-3} = 9\times 10^{-4}$،
و$\ln Q = -7.0$: فيكون $\Delta G =
-30.5 - 18.1 = \qty{-48.6}{kJ/mol}$.
وفي التعب: $Q = 3\times 10^{-3}\times
30\times 10^{-3}/5\times 10^{-3} = 0.018$،
و$\ln Q = -4.0$: فيكون $\Delta G =
-30.5 - 10.4 = \qty{-40.9}{kJ/mol}$.
فالتعب قصّ \qty{16}{\%} من الطاقة لكل \omterm{def:b1:water-small-molecules:atp}{ATP}، وهذا يكفي لإبطاء \omterm{def:b1:membranes-transport:transporters}{المضخات}
والمحركات التي تحتاج معظمها.
\end{solution}

\begin{solution}{exo:b1:water-small-molecules:11}
$[\mathrm{H^+}] = 10^{-7.2} = \qty{6.3e-8}{mol/L}$؛ والحجم
\qty{1e-12}{L}: أي $\qty{6.3e-20}{mol}\times\num{6e23} = 38$ بروتونًا.
مقابل $\num{3e9}$ زمرة \omterm{def:b1:water-small-molecules:buffer}{هادئة}: فليست «قيمة \omterm{def:b1:water-small-molecules:ph}{pH} في \omterm{def:b1:cell-unit-of-life:cell}{خلية}» عدًا للبروتونات بل حالة
برتنة الهوادئ، التي تتبادل البروتونات مع الماء أسرع بمليار مرة مما
يمكن للبروتونات الحرة القليلة أن تعنيه.
\end{solution}

\begin{solution}{exo:b1:water-small-molecules:12}
عند التوازن كان \omterm{def:b1:water-small-molecules:atp}{ATP} وADP والفوسفات ليستقروا عند نسبة تناهز
$10^{-5}$؛ لكن \omterm{def:b1:cell-unit-of-life:cell}{الخلية} تبقي \omterm{def:b1:water-small-molecules:atp}{ATP} فوق ADP بألف مرة، حتى تحرر حلمهته
\qty{50}{kJ/mol} بدل لا شيء. وكل عملية صعودية --- تركيب، ونقل، وحركة
--- مقترنة بذلك الانزياح؛ ودور الهدم الوحيد هو صونه. \omterm{def:b1:cell-unit-of-life:cell}{والخلية} التي
يبلغ \omterm{def:b1:water-small-molecules:atp}{ATP} فيها التوازن لا تدرج لها تنفقه: فتتوقف \omterm{def:b1:membranes-transport:transporters}{المضخات}، وترشح
الشوارد، وتنتفخ \omterm{def:b1:cell-unit-of-life:cell}{الخلية} وتموت في دقائق. فالحياة هي صون ذلك اللاتوازن،
لا مادة من المواد.
\end{solution}

\begin{solution}{pb:b1:water-small-molecules:1}
\textbf{1.} $6.1 + \log(24/1.2) = 6.1 + 1.301 = 7.40$.
\textbf{2.} $10^{-7.4} = \qty{4.0e-8}{mol/L} = \qty{40}{nmol/L}$.
\textbf{3.} $40\times 5 = \qty{200}{nmol}$، أي أقل من البيكربونات
بستمئة ألف مرة.
\textbf{4.} $\mathrm{CO_2} = 1.2\times 46/40 = \qty{1.38}{mmol/L}$؛
\omterm{def:b1:water-small-molecules:ph}{وpH} $= 6.1 + \log(25/1.38) = 6.1 + 1.258 = 7.36$.
\textbf{5.} في دورق لا تترك النسبة $20:1$ إلا قليلًا من الشريك
الحمضي لامتصاص \omterm{def:b1:water-small-molecules:ph}{الأساس}، والسعة منخفضة على بعد وحدة من p$K_a$؛ أما في
الجسم فالشريك الحمضي غاز تثبت الرئتان تركيزه وتستطيعان تجديده بلا
حد، فتُعاد النسبة إلى وضعها بالتهوية مهما كانت p$K_a$.
\textbf{6.} $24/1.2 = 20$؛ \omterm{def:b1:water-small-molecules:ph}{والأساس} $20/21 = \qty{95}{\%}$ من الزوج.
\textbf{7.} $60/5 = \qty{12}{mmol/L}$.
\textbf{8.} البيكربونات $24 - 12 = \qty{12}{mmol/L}$؛
و$\mathrm{CO_2}$ يصير $1.2 + 12 = \qty{13.2}{mmol/L}$.
\textbf{9.} \omterm{def:b1:water-small-molecules:ph}{pH} $= 6.1 + \log(12/13.2) = 6.1 - 0.04 = 6.06$.
\textbf{10.} $[\mathrm{H^+}] = \qty{0.012}{mol/L}$: أي \omterm{def:b1:water-small-molecules:ph}{pH} $1.92$.
\textbf{11.} إزاحة قدرها $1.34$ وحدة: فهي تنقذ الدم من \omterm{def:b1:water-small-molecules:ph}{pH} 1.9، لكن
6.06 أدنى بكثير مما يحتمله أي إنزيم. فهي رديئة، بوصفها جملة مغلقة.
\textbf{12.} ينتقل $[\mathrm{H^+}]$ من \qty{40}{nmol/L} إلى $10^{-6.06} = \qty{870}{nmol/L}$: أي ارتفاع قدره \qty{0.83}{\micro mol/L} مقابل \qty{12}{mmol/L} مضافة، أي بروتون واحد من كل \num{14000} يبقى حرًا؛ والباقي على البيكربونات.
\textbf{13.} \omterm{def:b1:water-small-molecules:ph}{pH} $= 6.1 + \log(12/1.2) = 6.1 + 1.0 = 7.10$.
\textbf{14.} $\mathrm{CO_2}$ المتولد:
$\qty{12}{mmol/L}\times
\qty{5}{L} = \qty{60}{mmol}$.
\textbf{15.} \qty{60}{mmol} في دقيقة فوق \qty{10}{mmol}: أي سبعة
أضعاف.
\textbf{16.} $\mathrm{CO_2} = \qty{0.9}{mmol/L}$؛ \omterm{def:b1:water-small-molecules:ph}{وpH} $= 6.1 +
\log(12/0.9) = 6.1 + 1.125 = 7.22$.
\textbf{17.} البيكربونات المفتوحة قرب 7.4: يتغير \omterm{def:b1:water-small-molecules:ph}{pH} بوحدة واحدة حين
تهبط البيكربونات عشرة أضعاف، من \qty{24}{} إلى \qty{2.4}{mmol/L}، أي
\qty{21.6}{mmol/L} من \omterm{def:b1:water-small-molecules:ph}{الحمض} لكل وحدة؛ لكن على الوحدات الأولى يكون
الميل الموضعي $2.3\times[\mathrm{HCO_3^-}] =
\qty{55}{mmol/L}$ لكل
وحدة. ومقابل \qty{25}{mmol/L} للبروتينات لكل وحدة، تأخذ البيكربونات
نحو $55/80 = \qty{69}{\%}$ من البروتونات، والبروتينات
\qty{31}{\%}.
\textbf{18.} تأخذ البيكربونات $0.69\times 12 = \qty{8.3}{mmol/L}$:
فتصير $24 - 8.3 = \qty{15.7}{mmol/L}$؛ \omterm{def:b1:water-small-molecules:ph}{وpH} $= 6.1 + \log(15.7/1.2) = 6.1 +
1.117 = 7.22$ (وحصة البروتينات متسقة مع \qty{25}{mmol/L} لكل
وحدة $\times 0.18$ وحدة $= \qty{4.5}{mmol/L}$، وهي قريبة من
\qty{3.7}{mmol/L}).
\textbf{19.} إزالة \omterm{def:b1:water-small-molecules:ph}{الحمض} تجدد البيكربونات المستهلكة
(\qty{12}{mmol/L} راجعة إلى \qty{24}{})، ومع تثبيت $\mathrm{CO_2}$
يعود \omterm{def:b1:water-small-molecules:ph}{pH} إلى 7.40.
\textbf{20.} $60/5 = \qty{12}{mmol/L}$ من البيكربونات تُفقد في
اليوم؛ وبعد ثلاثة أيام يكون $24 - 36 < 0$: فتنفد البيكربونات وينهار
\omterm{def:b1:water-small-molecules:ph}{pH} إلى ما دون 7 بكثير --- وعمليًا لا بد أن يُعطى المريض بيكربونات أو
يُغسل كلويًا؛ أما بعد يوم واحد وحده فيكون
$6.1 + \log(12/1.2) =
7.10$.
\textbf{21.} $\mathrm{CO_2} = 1.2\times 60/40 = \qty{1.8}{mmol/L}$؛
\omterm{def:b1:water-small-molecules:ph}{وpH} $= 6.1 + \log(24/1.8) = 6.1 + 1.125 = 7.22$. ومع بيكربونات عند
\qty{33}{}: $6.1 + \log(33/1.8) = 6.1 + 1.263 = 7.36$.
\textbf{22.} الرئتان تغيران المقام ($\mathrm{CO_2}$ المنحل) في
غضون أنفاس، لأنه غاز تزفرانه؛ والكليتان تغيران البسط (البيكربونات)
بطرح \omterm{def:b1:water-small-molecules:ph}{الحمض} وتجديد البيكربونات، وهي عملية ساعات إلى أيام.
\textbf{23.} عند \omterm{def:b1:water-small-molecules:ph}{pH} 5 يكون $[\mathrm{H^+}] = \qty{10}{\micro mol/L}$: أي \qty{15}{\micro mol} حرة في \qty{1.5}{L}، وهي جزء من أربعة آلاف من \qty{60}{mmol}؛ \omterm{def:b1:water-small-molecules:ph}{فالحمض} يغادر مرتبطًا بهوادئ --- أمونيومًا وفوسفاتًا ثنائي الهيدروجين.
\textbf{24.} فقد \omterm{def:b1:water-small-molecules:ph}{الحمض} يرفع \omterm{def:b1:water-small-molecules:ph}{pH}: فترتفع البيكربونات بمقدار
$100/5 = \qty{20}{mmol/L}$ إلى \qty{44}{}؛ \omterm{def:b1:water-small-molecules:ph}{وpH} $= 6.1 + \log(44/1.2) =
6.1 + 1.564 = 7.66$ (قلاء)، قبل أن تعوّض الرئتان والكليتان.
\textbf{25.} مغلقة: من 7.40 إلى 6.06، أي إزاحة $-1.34$. ومفتوحة
بالبيكربونات وحدها: إلى 7.10، أي إزاحة $-0.30$. ومفتوحة مع
البروتينات: إلى 7.22، أي إزاحة $-0.18$.
\end{solution}
